\documentclass[conference]{IEEEtran}
\usepackage{amsmath}
\usepackage{booktabs}
\usepackage{cite}
\author{\IEEEauthorblockN{Ritik, Aman, Ansh Tomar, Sarabjeet Singh, and Deepanshu Chaudhary}}


\title{A Two-Way LRU Cache Prototype in Verilog\\
\large Phase I Project Report}
\begin{document}
\maketitle

\begin{abstract}
We developed a Verilog prototype to show how cache associativity changes
block placement and replacement. The prototype pairs a direct-mapped baseline
with a two-way cache that uses least recently used (LRU) replacement.
Both designs hold 64 bytes in four-byte blocks. A directed test fills both
ways of one set, reads the older block again, and then requests a third
block. The two-way controller keeps the block that was just read and evicts
the other one. A shared testbench checks this behavior, returned bytes,
and time to complete each access.
\end{abstract}

\begin{IEEEkeywords}
Cache memory, Verilog, set associativity, least recently used
\end{IEEEkeywords}

\section{Introduction}
A cache retains recently accessed blocks so the CPU can satisfy some
requests without waiting for main memory. A direct-mapped cache gives a
block one possible location. Two frequently accessed blocks that share
that location can displace each other even when other lines are empty.
The design question for this prototype is whether two slots in a selected
set retain such blocks at the same total data capacity. The prototype also
makes the victim selection visible when a third block needs the same set.

\section{Literature Review}
Smith surveys cache design choices such as placement, replacement, block size,
and memory update policy \cite{smith}. Placement limits the possible slots
for a block. Replacement chooses a slot when those locations are occupied.
Hill and Smith use address traces to study how associativity changes miss
ratios and to isolate causes of misses \cite{hill}. We use short, identical
request sequences to show exactly which address causes each hit, miss,
or eviction in the prototype.

\section{Proposed Solution}
The testbench supplies byte reads and writes to one cache at a time. The
controller saves the request when idle, then compares tags in the selected
location. A hit requires a matching tag and an asserted valid bit. The two-way
design compares its two tags in parallel. On a miss, the controller fetches
the aligned four-byte block from the memory model.

\begin{table}[!t]
\caption{Cache Geometry and Address Fields}
\label{tab:geom1}
\centering
\begin{tabular}{lcc}
\toprule
 & Direct-mapped & Two-way LRU\\
\midrule
Sets $\times$ ways & $16\times1$ & $8\times2$\\
Data capacity & 64 B & 64 B\\
Block size & 4 B & 4 B\\
Tag bits & $[15:6]$ & $[15:5]$\\
Index bits & $[5:2]$ & $[4:2]$\\
Byte offset bits & $[1:0]$ & $[1:0]$\\
\bottomrule
\end{tabular}
\end{table}

\subsection{Decoding an Address}
Table~\ref{tab:geom1} defines the 16-bit byte address fields. For example,
\texttt{0x0043} selects two-way set 0 with tag 2. Its offset is 3, so it
uses the fourth byte of the block starting at \texttt{0x0040}. The same
address maps to direct line 0 with tag 1.

\subsection{Filling and Replacing Blocks}
A two-way miss uses an invalid slot first. When both ways are valid, the
set's LRU bit points to the way to replace. A hit or fill makes its way
most recently used, so the controller saves the other way as the next
victim. One bit suffices because the set holds just two ways.

\subsection{Finishing Reads and Writes}
A read returns one byte. A write updates only the addressed byte. With
write-allocate, a write miss first fetches the block. With write-through,
every write then sends the selected byte to memory and waits for
acknowledgment. The seven controller states accept the CPU request, check
tags, issue and wait for a refill, issue and wait for a write, and return
the response. Read hits proceed from lookup directly to response.

\section{Prototype Verification}
Table~\ref{tab:lru} shows the first four requests of the directed demo.
The third request hits A and makes B the next victim. C then evicts B.
The testbench follows up by reading A again and then B. A hits and B
misses, confirming the contents of the set through the normal CPU interface.
Each block maps to line 0 in direct mapping, where each different block
overwrites the previous one. Across these six reads, direct mapping
takes 42 service cycles with zero hits and the two-way cache takes
32 cycles with two hits.

\begin{table}[!t]
\caption{Two-Way Set 0 After Each Demo Access}
\label{tab:lru}
\centering
\footnotesize
\begin{tabular}{lcccl}
\toprule
Read address & Result & Way 0 & Way 1 & Next victim\\
\midrule
A (\texttt{0000}) & Miss & A & Empty & Empty way\\
B (\texttt{0040}) & Miss & A & B & A\\
A (\texttt{0000}) & Hit & A & B & B\\
C (\texttt{0080}) & Miss & A & C & A\\
\bottomrule
\end{tabular}
\end{table}

\subsection{Reproducing the Demonstration}
From the project root, \texttt{./run.sh demo} compiles both designs with
Icarus Verilog, prints each result, and saves a waveform for each cache.
In \texttt{2way\_demo.vcd}, \texttt{cpu\_req} marks the request
and \texttt{cpu\_done} marks completion. On a miss, \texttt{mem\_req}
requests the block and \texttt{mem\_done} acknowledges it. Figure~\ref{fig:wave}
shows these signals on the third and fourth reads. Counting from the
accepting clock edge, the hit completes at 225 ns (two cycles) and the
miss completes at 305 ns (seven cycles). The final two reads then
confirm the eviction through the CPU interface.

\begin{figure}[!t]
\centering
\setlength{\unitlength}{1pt}
\begin{picture}(236,116)
\scriptsize
\put(85,105){\makebox(0,0){A hits}}
\put(163,105){\makebox(0,0){C misses}}
\put(5,88){\texttt{cpu\_req}}
\put(5,69){\texttt{mem\_req}}
\put(5,50){\texttt{mem\_done}}
\put(5,31){\texttt{cpu\_done}}
\thicklines
\put(63,93){\line(1,0){15}}
\put(78,86){\line(0,1){7}}
\put(78,86){\line(1,0){30}}
\put(108,86){\line(0,1){7}}
\put(108,93){\line(1,0){15}}
\put(123,86){\line(0,1){7}}
\put(123,86){\line(1,0){105}}
\put(63,67){\line(1,0){67.5}}
\put(130.5,67){\line(0,1){7}}
\put(130.5,74){\line(1,0){15}}
\put(145.5,67){\line(0,1){7}}
\put(145.5,67){\line(1,0){82.5}}
\put(63,48){\line(1,0){127.5}}
\put(190.5,48){\line(0,1){7}}
\put(190.5,55){\line(1,0){15}}
\put(205.5,48){\line(0,1){7}}
\put(205.5,48){\line(1,0){22.5}}
\put(63,36){\line(1,0){7.5}}
\put(70.5,29){\line(0,1){7}}
\put(70.5,29){\line(1,0){30}}
\put(100.5,29){\line(0,1){7}}
\put(100.5,36){\line(1,0){15}}
\put(115.5,29){\line(0,1){7}}
\put(115.5,29){\line(1,0){105}}
\put(220.5,29){\line(0,1){7}}
\put(220.5,36){\line(1,0){7.5}}
\thinlines
\put(63,17){\line(1,0){165}}
\multiput(63,14)(30,0){6}{\line(0,1){3}}
\put(63,7){\makebox(0,0){200}}
\put(93,7){\makebox(0,0){220}}
\put(123,7){\makebox(0,0){240}}
\put(153,7){\makebox(0,0){260}}
\put(183,7){\makebox(0,0){280}}
\put(213,7){\makebox(0,0){300}}
\put(228,7){\makebox(0,0){ns}}
\end{picture}
\caption{Signals from the two-way demo for the second read of A
and the first read of C. The testbench clock period is 10 ns.}
\label{fig:wave}
\end{figure}


\end{document}
